% ------------------------------------------------------------------------------
% 4 Results
% ------------------------------------------------------------------------------
\section{Results}\label{sec:results}

\subsection{Treatment effects on disease-free and overall survival}
\label{sec:res-tx}

\subsubsection*{Kaplan--Meier analysis}

Figures~\ref{fig:km} and~\ref{fig:kmos} display Kaplan--Meier estimates
for the composite DFS endpoint and for overall survival, each with the
numbers at risk. The three curves separate early and the separation
widens over the first three years, driven almost entirely by the
levamisole-plus-5-FU arm. The overall log-rank test is highly
significant for DFS ($\chi^2 = 21.0$ on 2 degrees of freedom,
$p < 0.001$); pairwise comparisons show that levamisole alone is
indistinguishable from observation (HR $0.968$, $p = 0.76$), whereas
levamisole plus 5-FU is superior to both observation (HR $0.623$,
$p < 0.001$) and to levamisole alone ($p < 0.001$). The pattern for
overall survival is similar but attenuated (adjusted HR $0.688$, 95\% CI
$0.544$--$0.871$), consistent with post-recurrence salvage therapies
diluting survival differences.

\begin{figure}[htbp]
  \centering
  \includegraphics[width=\textwidth]{fig_km_dfs.pdf}
  \caption{Kaplan--Meier estimates of the composite disease-free
  survival endpoint (primary) with pointwise 95\% confidence bands and
  numbers at risk.}
  \label{fig:km}
\end{figure}

\begin{figure}[htbp]
  \centering
  \includegraphics[width=\textwidth]{fig_km_os.pdf}
  \caption{Kaplan--Meier estimates of overall survival (secondary
  endpoint) with pointwise 95\% confidence bands and numbers at risk.}
  \label{fig:kmos}
\end{figure}

\subsubsection*{Cox models}

Table~\ref{tab:cox} assembles the treatment-effect estimates across
specifications. Adjustment for the ten baseline covariates, pooled over
20 multiple imputations by Rubin's rules, moves the primary contrast
only marginally (unadjusted HR $0.623$ vs.\ adjusted $0.621$), as
expected under randomization with balanced arms; the complete-case
analysis agrees closely. The adjusted covariate effects themselves
(Figure~\ref{fig:forest}) are clinically coherent: each additional
positive lymph node multiplies the hazard by $1.08$ (95\% CI
$1.06$--$1.10$), spread to contiguous structures triples the hazard
relative to submucosal involvement (HR $3.19$, 95\% CI $1.32$--$7.75$),
obstruction raises it by 25\%, and a long surgery-to-registration
interval is protective in this cohort (HR $0.76$, 95\% CI $0.62$--$0.92$).

\begin{table}[htbp]
\centering\small
\caption{Treatment-effect estimates for the composite DFS endpoint
across analysis specifications. The Fine--Gray model targets the
recurrence endpoint with death without recurrence as a competing event.}
\label{tab:cox}
\begin{tabular*}{\textwidth}{@{\extracolsep{\fill}} l l c c c}
\toprule
Analysis & Contrast & HR & 95\% CI & $p$ \\
\midrule
Unadjusted Cox & Lev vs.\ Obs & 0.968 & 0.790--1.187 & 0.757 \\
Unadjusted Cox & Lev+5FU vs.\ Obs & 0.623 & 0.499--0.777 & $<0.001$ \\
Adjusted Cox, MI ($m=20$) & Lev vs.\ Obs & 0.947 & 0.770--1.165 & 0.607 \\
Adjusted Cox, MI ($m=20$) & Lev+5FU vs.\ Obs & \textbf{0.621} & \textbf{0.496--0.776} & $<0.001$ \\
Adjusted Cox, MI (OS endpoint) & Lev+5FU vs.\ Obs & 0.688 & 0.544--0.871 & 0.002 \\
Fine--Gray subdistribution & Lev+5FU vs.\ Obs & 0.596 & 0.469--0.756 & $<0.001$ \\
\bottomrule
\end{tabular*}
\end{table}

\begin{figure}[htbp]
  \centering
  \includegraphics[width=0.95\textwidth]{fig_forest_adjusted.pdf}
  \caption{Adjusted hazard ratios for disease-free survival from the
  multiply imputed Cox model (log scale). Red indicates estimates whose
  95\% confidence interval excludes 1.}
  \label{fig:forest}
\end{figure}

\subsubsection*{Absolute benefit}

Five-year Kaplan--Meier DFS probabilities were 42.4\% (95\% CI
36.9--47.8) under observation, 44.2\% (38.6--49.6) under levamisole,
and 59.2\% (53.4--64.5) under levamisole plus 5-FU; the corresponding
overall-survival probabilities were 52.6\%, 53.5\%, and 63.4\%. The
bootstrap 95\% CI for the five-year DFS risk difference of levamisole
plus 5-FU against observation is $-16.7$ percentage points ($-24.2$ to
$-9.4$), corresponding to a number needed to treat of six patients to
prevent one additional DFS event at five years. Levamisole alone shows
no detectable absolute benefit ($-1.8$ points, $-9.3$ to $+5.7$).

\subsubsection*{Proportional-hazards diagnostics}

The scaled Schoenfeld-residual analysis flags two covariates: tumour
differentiation ($p < 0.001$) and obstruction ($p = 0.008$), with a
significant global test ($\chi^2 = 42.2$, 14 df, $p < 0.001$).
Reassuringly, the treatment indicator itself shows no evidence of a
time-varying effect ($p = 0.51$), so the primary hazard-ratio summary is
defensible. Figure~\ref{fig:zph} shows the residual patterns; the
obstruction effect visibly weakens after year two. Replacing linear age
and nodes with 4-degree-of-freedom penalized splines significantly
improves the log-likelihood (likelihood-ratio $p < 0.001$) without
affecting the treatment estimate.

\begin{figure}[htbp]
  \centering
  \includegraphics[width=\textwidth]{fig_zph.pdf}
  \caption{Proportional-hazards diagnostics: scaled Schoenfeld
  residuals with loess smooths for treatment, positive nodes, and
  obstruction (which exhibits a significant violation, $p = 0.008$).}
  \label{fig:zph}
\end{figure}

\subsubsection*{Competing-risk sensitivity analysis}

Thirty-eight patients died without a recorded recurrence, and treating
these deaths as censoring---as the published coding does---formally
assumes they carry no recurrence-risk information. The competing-risk
view shows the practical consequence (Figure~\ref{fig:cif}): the
Aalen--Johansen cumulative incidence of recurrence sits below the naive
one-minus-Kaplan--Meier curve for every arm, with the gap reaching about
one percentage point at eight years in the observation arm. The
Fine--Gray subdistribution model estimates the levamisole-plus-5-FU
effect on the recurrence subdistribution hazard at $0.596$ (95\% CI
$0.469$--$0.756$)---virtually identical to the cause-specific reading of
$0.623$---so the conclusion of a substantial treatment benefit is robust
to the endpoint-definition choice.

\begin{figure}[htbp]
  \centering
  \includegraphics[width=0.95\textwidth]{fig_cif.pdf}
  \caption{Cumulative incidence of recurrence: Aalen--Johansen
  estimates treating death without recurrence as a competing event
  (solid) versus the naive one-minus-Kaplan--Meier estimator that
  censors such deaths (dashed).}
  \label{fig:cif}
\end{figure}

\subsubsection*{Restricted mean survival time}

The RMST analysis restates the benefit in time units without any
proportional-hazards assumption (Table~\ref{tab:rmst},
Figure~\ref{fig:rmst}). Within a five-year horizon, patients randomized
to levamisole plus 5-FU can expect 3.57 disease-free years on average
versus 2.94 under observation---a gain of 0.63 years (95\% CI 0.33 to
0.92), roughly seven and a half months, with an RMST ratio of 1.21
(1.11 to 1.33). The gain widens with the horizon, reaching 1.13 years
(0.62 to 1.65) by eight years as the curves continue to separate. The
covariate-adjusted pseudo-observation regression yields 0.57 years
(0.28 to 0.85), closely agreeing with the unadjusted estimate---as it
must, because the RMST difference, unlike the hazard ratio, is
collapsible: covariate adjustment changes precision, not the estimand.
Within the obstruction stratum---the covariate with the strongest
proportional-hazards violation ($p = 0.008$)---the five-year RMST
difference remains substantial (0.60 years, 0.27 to 0.94, without
obstruction; 0.83 years, 0.11 to 1.55, with obstruction), confirming
that the treatment conclusion does not rest on the proportional-hazards
assumption.

\begin{table}[htbp]
\centering\small
\caption{Restricted mean survival time for the primary comparison
(Lev+5FU vs.\ observation, DFS endpoint). Difference $=$ disease-free
years gained; percentile bootstrap CIs (2{,}000 resamples).}
\label{tab:rmst}
\begin{tabular*}{\textwidth}{@{\extracolsep{\fill}} c c c c c}
\toprule
Horizon $\tau$ & RMST Obs (y) & RMST Lev+5FU (y) & Difference (95\% CI) & Ratio (95\% CI) \\
\midrule
3 & 2.03 & 2.34 & 0.31 (0.14--0.47) & 1.15 (1.07--1.24) \\
4 & 2.50 & 2.96 & 0.46 (0.23--0.68) & 1.18 (1.09--1.29) \\
5 & 2.94 & 3.57 & \textbf{0.63 (0.33--0.92)} & 1.21 (1.11--1.33) \\
6 & 3.36 & 4.14 & 0.79 (0.42--1.15) & 1.23 (1.12--1.37) \\
7 & 3.75 & 4.70 & 0.95 (0.51--1.39) & 1.25 (1.13--1.40) \\
8 & 4.12 & 5.24 & 1.13 (0.62--1.65) & 1.27 (1.14--1.43) \\
\bottomrule
\end{tabular*}
\end{table}

\begin{figure}[htbp]
  \centering
  \includegraphics[width=0.95\textwidth]{fig_rmst.pdf}
  \caption{Restricted mean survival time analysis. Left: the RMST is the
  shaded area under the Kaplan--Meier curve up to the horizon $\tau = 5$
  years (shown for the observation arm). Right: RMST difference
  (disease-free years gained with levamisole plus 5-FU) with 95\%
  bootstrap confidence intervals as a function of the horizon.}
  \label{fig:rmst}
\end{figure}

\subsection{Subgroups, trial redesign, and re-enacted monitoring}
\label{sec:res-design}

\subsubsection*{Subgroup analyses}

Figure~\ref{fig:subgroup} presents the seven pre-specified subgroup
analyses for the primary contrast. The hazard ratio is remarkably stable
across subgroups---every point estimate lies between $0.46$ and $0.81$
and every confidence interval overlaps the overall estimate of $0.62$.
Only one nominal interaction reaches $p<0.05$: sex (interaction
$p = 0.043$), with a substantially larger effect in men (HR $0.46$,
95\% CI $0.33$--$0.64$) than in women (HR $0.79$, 95\% CI $0.58$--$1.07$).
Against the Bonferroni-adjusted threshold of $0.0071$ appropriate for
seven tests, this signal does not survive, and no biological mechanism
has been established for sex-specific efficacy of 5-FU-based regimens in
this setting. The honest reading is \emph{consistent relative effect
across all examined subgroups}.

\begin{figure}[htbp]
  \centering
  \includegraphics[width=\textwidth]{fig_subgroup_forest.pdf}
  \caption{Pre-specified subgroup analyses for disease-free survival,
  levamisole plus 5-FU versus observation. Right column: hazard ratios
  with 95\% CIs and interaction-test $p$-values (adjusted for nodal
  count).}
  \label{fig:subgroup}
\end{figure}

\subsubsection*{The redesigned trial}

Table~\ref{tab:design} summarizes the fixed design and the
group-sequential redesign. For the target hazard ratio of $0.65$ at
one-sided $\alpha = 0.025$ with 90\% power, the fixed design requires
226 DFS events (approximately 322 patients under the assumed accrual and
follow-up). The three-look sequential design inflates the maximum to 235
events but compensates with expected savings whenever the effect is real
or absent. The O'Brien--Fleming efficacy boundary is naturally
conservative early ($Z = 2.963$ at the first interim) and relaxes toward
$2.014$ at the final analysis; the non-binding futility boundary starts
at $Z = 0.079$, meaning a first look with essentially no observed trend
would trigger the futility signal. Power falls steeply for weaker
effects: at $\hr = 0.80$, nearly 850 events would be required---a useful
reminder of why ambitious trials with modest expected effects grow so
large.

\begin{table}[htbp]
\centering\footnotesize
\caption{The redesigned group-sequential trial: boundaries and operating
characteristics from \texttt{gsDesign}. Efficacy uses O'Brien--Fleming
alpha-spending; futility uses a non-binding Hwang--Shih--DeCani
beta-spending ($\gamma = -4$).}
\label{tab:design}
\begin{tabular*}{\textwidth}{@{\extracolsep{\fill}} l c c c c c c}
\toprule
Analysis & Events & Effic.\ $Z$ & 1-sided $p$ & HR at bound & Futil.\ $Z$ & HR at bound \\
\midrule
Interim 1 (50\%)  & 118 & 2.963 & 0.0015 & 0.579 & 0.078 & 0.986 \\
Interim 2 (75\%)  & 177 & 2.359 & 0.0091 & 0.701 & 1.006 & 0.859 \\
Final (100\%)     & 235 & 2.014 & 0.0220 & 0.769 & 2.014 & 0.769 \\
\midrule
\multicolumn{7}{@{}l@{}}{\emph{Design-theoretic probabilities:}
$P(\text{reject} \mid \mathrm{HR}=0.65) = 0.900$;\;
$P(\text{reject} \mid \mathrm{HR}=1) = 0.024$ (futility adhered)} \\
\bottomrule
\end{tabular*}
\end{table}

\begin{figure}[htbp]
  \centering
  \includegraphics[width=0.95\textwidth]{fig_gs_boundaries.pdf}
  \caption{Group-sequential boundaries of the redesigned trial on the
  $Z$ scale (left axis) and one-sided $p$-value scale (right axis).
  Red: efficacy (O'Brien--Fleming spending); blue: non-binding
  futility (Hwang--Shih--DeCani).}
  \label{fig:gsbound}
\end{figure}

\subsubsection*{Re-enacting interim monitoring on the real data}

The most striking result of this re-analysis emerges when the redesigned
monitoring plan is applied to the \emph{actual} event stream. Censoring
every patient's follow-up at the time at which the 118th DFS event was
observed (information fraction 50\%) and refitting the two-arm Cox model
yields an interim hazard ratio of $0.52$ and a $Z$-statistic of
$3.38$---\emph{above the efficacy boundary of $2.963$}
(Figure~\ref{fig:reenact}). Under the redesign, the trial would have
stopped at the first interim analysis and declared the levamisole plus
5-FU regimen superior to observation, after only 118 of the 506 events
that were eventually observed. Conditional power calculations quantify
how decisive this interim result is: the probability of eventually
crossing the final boundary is 70.6\% even under the null, 99.8\% under
the design alternative, and essentially 100\% if the observed interim
trend persists. A cross-check with \texttt{gsDesign}'s conditional-power
machinery returns somewhat higher values ($0.834$ and $0.999$ under the
null and alternative), because it accumulates the crossing probabilities
at \emph{both} remaining analyses with boundaries re-optimized after the
interim, whereas the hand-rolled calculation uses a single-look normal
approximation to the final analysis only. The contrast with the Bayesian
view is instructive: sequential posterior probabilities under the
skeptical prior (Section~\ref{sec:res-bayes}) would have signalled
benefit even earlier, at the 25\% information look.

Two interpretive cautions apply. First, the historical trial was not
designed with this monitoring plan, so the re-enactment is a
counterfactual exercise, not a critique of the original conduct; its
message is about how much evidence was already present early in the
event stream. Second, the interim hazard ratio of $0.52$ is more extreme
than the final estimate of $0.62$, a reminder that early random highs in
$Z$ are partly regression artefacts; the design's conservative early
boundary exists precisely to price this randomness in.

\begin{figure}[htbp]
  \centering
  \includegraphics[width=0.95\textwidth]{fig_gs_reenactment.pdf}
  \caption{Re-enacted interim monitoring using the real trial data.
  The observed interim result after 118 events (HR $0.52$, $Z = 3.38$,
  green marker) crosses the O'Brien--Fleming efficacy boundary
  ($2.96$), so the redesigned trial would have stopped for efficacy at
  the first interim analysis.}
  \label{fig:reenact}
\end{figure}

\subsection{Simulation studies}\label{sec:res-sim}

\subsubsection*{Study A: covariate adjustment buys power}

Table~\ref{tab:simA} and Figure~\ref{fig:simA} report 1{,}000 simulated
trials per scenario at $n = 322$ and target hazard ratio $0.72$. With no
prognostic covariates, adjustment neither helps nor hurts, as theory
predicts: power is 71.6\% unadjusted versus 70.4\% adjusted, and the
type-I error is 2.5\% and 2.9\% respectively (2.2\% and 3.2\% at full
strength)---all within Monte Carlo error of the nominal 2.5\%. As
prognostic strength grows to the trial-calibrated value (SD of the
prognostic index $0.416$ on the log-hazard scale), adjusted power rises
to 80.9\% against 72.6\% unadjusted---an efficiency gain equivalent to
roughly 70 additional events (about a 23\% larger event count), obtained
at zero cost in bias or type-I error. This is the quantitative case
behind current regulatory encouragement of pre-specified covariate
adjustment \citep{fda2023,kahan2014}.

The bias column also teaches a subtler lesson: at full prognostic
strength, the unadjusted estimator averages $+0.033$ on the log scale
relative to the data-generating conditional hazard ratio---that is, an
unadjusted marginal HR of about $0.745$ when the conditional truth is
$0.72$. This attenuation toward the null is the hazard ratio's
non-collapsibility \citep{martinussen2013}: unlike a linear-model
coefficient or a risk difference (or, as demonstrated above, the RMST
difference), the marginal HR changes when strong prognostic covariates
are removed from the model even without any confounding. It explains why
adjusted and unadjusted estimates in randomized trials need not
coincide, and why the near-coincidence of $0.623$ and $0.621$ in our
real data reflects the modest prognostic strength of the adjustment set
rather than a mathematical necessity.

\begin{table}[htbp]
\centering\small
\caption{Simulation Study A: covariate-adjustment efficiency (1{,}000
replicates per scenario, $n=322$, one-sided $\alpha=0.025$). Power at
HR $=0.72$; type-I error at HR $=1$. ``Bias'' contrasts each estimator
with the conditional log-HR used to generate the data.}
\label{tab:simA}
\begin{tabular*}{\textwidth}{@{\extracolsep{\fill}} l l c c c c}
\toprule
Prognostic strength & HR & Power (unadj.) & Power (adj.) & Bias (unadj.) & Bias (adj.) \\
\midrule
None (SD $0.000$)     & 1.00 & 0.025 & 0.029 & $-0.006$ & $-0.006$ \\
Moderate (SD $0.208$) & 1.00 & 0.025 & 0.025 & $+0.004$ & $+0.005$ \\
Full (SD $0.416$)     & 1.00 & 0.022 & 0.032 & $+0.000$ & $-0.001$ \\
None (SD $0.000$)     & 0.72 & 0.716 & 0.704 & $+0.000$ & $-0.009$ \\
Moderate (SD $0.208$) & 0.72 & 0.725 & 0.766 & $+0.012$ & $-0.006$ \\
Full (SD $0.416$)     & 0.72 & 0.726 & \textbf{0.809} & $+0.033$ & $-0.010$ \\
\bottomrule
\end{tabular*}
\end{table}

\begin{figure}[htbp]
  \centering
  \includegraphics[width=0.95\textwidth]{fig_sim_efficiency.pdf}
  \caption{Simulation Study A: empirical power of the unadjusted versus
  covariate-adjusted Cox test as a function of prognostic strength
  (target HR $0.72$, $n = 322$, 1{,}000 replicates per scenario).}
  \label{fig:simA}
\end{figure}

\subsubsection*{Study B: operating characteristics of the sequential
design}

Simulating the full monitoring procedure (3{,}000 trials per scenario,
event-triggered looks with calendar-time censoring) validates the design
and quantifies the value of the futility boundary
(Table~\ref{tab:simB}, Figure~\ref{fig:simB}). Under the alternative
(HR $0.65$), the design rejects in 82.9\% of trials against the
design-theoretic 90\%. Three mechanisms account for the shortfall, all
instructive. First, the simulated primary analysis is the
\emph{unadjusted} Cox test, whose target is the marginal hazard
ratio---attenuated toward the null by the strong calibration covariates
through exactly the non-collapsibility quantified in Study A. Second,
futility adherence stops 6.8\% of trials under the alternative that
might eventually have crossed. Third, event-triggered looks with
discrete calendar-time cuts approximate, but do not reproduce, the
continuous information accrual assumed by the spending functions.
Stopping is dominated by the second interim (42.2\%) rather than the
first (18.6\%), and the expected event count is 183.4 versus 235 at
completion: the design saves, on average, 51.6 events (roughly 22\%)
when the drug works. Under the null with futility adhered to, the type-I
error is 2.4\% and the trial stops for futility at the first interim in
54.9\% of cases. Ignoring futility (the non-binding option) inflates the
type-I error only trivially, to 2.7\%, confirming the design's stated
robustness.

\begin{table}[htbp]
\centering\footnotesize
\caption{Simulation Study B: simulated operating characteristics of the
group-sequential design (3{,}000 trials per scenario).}
\label{tab:simB}
\begin{tabular*}{\textwidth}{@{\extracolsep{\fill}} l c c c c c c}
\toprule
Scenario & Eff.\ IA1 & Eff.\ IA2 & Final & Total eff. & Exp.\ events & Events saved \\
\midrule
HR $=0.65$ (as designed)   & 0.186 & 0.422 & 0.222 & \textbf{0.829} & 183.4 & 51.6 \\
HR $=1$, futility adhered  & 0.002 & 0.008 & 0.014 & 0.024 & 152.7 & 82.4 \\
HR $=1$, futility ignored  & 0.001 & 0.009 & 0.017 & 0.027 & 234.4 & 0.6 \\
\bottomrule
\end{tabular*}
\end{table}

\begin{figure}[htbp]
  \centering
  \includegraphics[width=\textwidth]{fig_sim_gs_oc.pdf}
  \caption{Simulation Study B: simulated stopping probabilities by
  analysis and reason under the alternative (right) and under the null
  with futility adherence (left).}
  \label{fig:simB}
\end{figure}

\subsubsection*{Study C: interval estimation under non-proportional
hazards}

When the true effect is time-varying (HR $0.45$ in year one, weakening
to $0.80$ thereafter) and the analysis fits a single
proportional-hazards Cox model, the estimator targets a weighted average
of the two regimes (simulation estimand: log-HR $= -0.407$, i.e.\ a
pooled HR of about $0.666$). Over 500 replicates, the Wald interval
achieved 92.6\% coverage of this estimand and the percentile bootstrap
90.4\%---both marginally below their nominal 95\%
(Table~\ref{tab:simC}). The bootstrap is slightly narrower (mean width
$0.434$ vs.\ $0.443$) but not more accurate. The lesson is double-edged:
neither standard interval collapses under moderate non-proportional
hazards---the practical damage is a shifted estimand rather than
drastically wrong uncertainty---yet the modest undercoverage confirms
that misspecification costs are real, reinforcing the value of the
diagnostics and of the RMST analysis of Section~\ref{sec:res-tx}.

\begin{table}[htbp]
\centering\footnotesize
\caption{Simulation Study C: 95\% interval performance under a
time-varying hazard ratio (500 replicates, $n=335$; 100 bootstrap
resamples per replicate).}
\label{tab:simC}
\begin{tabular*}{0.80\textwidth}{@{\extracolsep{\fill}} l c c c}
\toprule
Method & Coverage & Mean width & Mean log-HR \\
\midrule
Cox Wald CI            & 0.926 & 0.443 & $-0.407$ \\
Bootstrap percentile & 0.904 & 0.434 & $-0.407$ \\
\bottomrule
\end{tabular*}
\end{table}

\begin{figure}[htbp]
  \centering
  \includegraphics[width=0.95\textwidth]{fig_sim_coverage.pdf}
  \caption{Simulation Study C: Wald versus bootstrap percentile
  intervals for the log hazard ratio across 500 simulated trials under
  a time-varying true effect. The dashed line marks the simulation
  estimand.}
  \label{fig:simC}
\end{figure}

\subsection{Machine-learning prediction and Bayesian inference}
\label{sec:res-ml}

\subsubsection*{Cross-validated prediction performance}

Table~\ref{tab:ml} and Figure~\ref{fig:mlcv} compare the five prognostic
models over 25 cross-validation folds (888 complete cases). The random
survival forest achieves the best discrimination (cross-validated
$C = 0.646$, IPCW AUC at five years $0.677$) and the lowest integrated
Brier score ($0.217$ over years 1--5), ahead of the clinical Cox model
($C = 0.634$), gradient boosting ($0.633$), XGBoost ($0.625$), and the
elastic-net Cox model ($0.620$). The spread among models is
modest---about 0.026 in the $C$-index---and the honest uncertainty
(fold-to-fold SD $\approx 0.025$) cautions against declaring a
definitive winner; what the comparison really shows is that, with ten
baseline covariates and $\sim$300 events, flexible learners extract only
a little more signal than a well-specified clinical model.

The apparent-versus-cross-validated contrast in Table~\ref{tab:ml}
quantifies overfitting: XGBoost's apparent $C$ of $0.741$ drops to
$0.625$ under honest validation, a 0.12 optimism gap, versus a 0.07 gap
for the random forest and essentially none for the clinical Cox model.
This gap is the single most important column in the table---it is the
difference between the performance one would report from a naive
in-sample fit and the performance a new patient would actually
experience. Calibration of the winning model
(Figure~\ref{fig:mlcal}) is satisfactory across the risk spectrum:
within deciles of predicted five-year risk, the mean predicted risk
tracks the Kaplan--Meier observed risk within $6.5$ percentage points.
Permutation importance is dominated by the positive lymph-node count
(mean $C$-index drop $0.117$ per permutation, an order of magnitude
above every other covariate), followed at a distance by advanced local
extent, female sex, and a long surgery-to-registration interval---a
ranking that mirrors the adjusted Cox effects of Figure~\ref{fig:forest}
and reinforces the clinical coherence of the learned model.

\begin{table}[htbp]
\centering\small
\caption{Machine-learning prediction of the composite DFS endpoint
(888 complete cases; $5\times5$-fold stratified cross-validation).
Treatment excluded from all predictors. AUC at five years is IPCW
weighted; IBS is the integrated Brier score over years 1--5.}
\label{tab:ml}
\begin{tabular*}{\textwidth}{@{\extracolsep{\fill}} l c c c c}
\toprule
Model & CV $C$-index & CV AUC (5y) & CV IBS & Apparent $C$ \\
\midrule
Random survival forest         & \textbf{0.646} & \textbf{0.677} & \textbf{0.217} & 0.713 \\
Clinical Cox (main effects)    & 0.634 & 0.652 & 0.222 & 0.651 \\
Gradient boosting (Cox loss)   & 0.633 & 0.662 & 0.225 & 0.717 \\
XGBoost (Cox objective) & 0.625 & -- & -- & 0.741 \\
Elastic-net Cox                & 0.620 & 0.646 & 0.226 & 0.628 \\
\bottomrule
\end{tabular*}
\end{table}

\begin{figure}[htbp]
  \centering
  \includegraphics[width=0.95\textwidth]{fig_ml_cv_performance.pdf}
  \caption{Cross-validated discrimination by model. Left: distribution
  of the $C$-index across 25 folds. Right: IPCW AUC at five years with
  cross-validation standard errors.}
  \label{fig:mlcv}
\end{figure}

\begin{figure}[htbp]
  \centering
  \includegraphics[width=0.95\textwidth]{fig_ml_calibration_auc.pdf}
  \caption{Discrimination and calibration of the random survival
  forest. Left: ROC curve for five-year event classification. Right:
  predicted versus Kaplan--Meier observed five-year event risk by decile
  of predicted risk.}
  \label{fig:mlcal}
\end{figure}

\subsubsection*{Bayesian re-analysis: priors and posteriors}
\label{sec:res-bayes}

The Weibull model converged cleanly under all three prior scenarios
(two chains, target acceptance $0.9$; $\hat R \approx 1.00$ throughout).
Table~\ref{tab:bayes} and Figure~\ref{fig:bayesprior} display the
central result: the posterior hazard ratio is essentially insensitive to
the prior, ranging only from $0.591$ (weakly-informative) to $0.610$
(skeptical). Under the skeptical prior---which places 95\% of its mass
between hazard ratios of about $0.44$ and $2.28$ and only 5\% below
$0.5$---the posterior mean is $0.610$ with 95\% credible interval
$0.485$--$0.757$, $\Prob(\mathrm{HR}<1) = 1.000$ to three decimals, and
$\Prob(\mathrm{HR}<0.8) = 0.994$. The data are simply strong enough to
overwhelm any of the three priors, which is itself an informative
finding: conclusions of substance do not require prior consensus.

\begin{table}[htbp]
\centering\footnotesize
\caption{Bayesian Weibull analysis of the primary contrast under three
prior specifications (2 chains; 4{,}000 posterior draws each), with the
frequentist adjusted estimate for comparison.}
\label{tab:bayes}
\begin{tabular*}{\textwidth}{@{\extracolsep{\fill}} l c c c c c}
\toprule
Prior & Mean (log-HR) & Post.\ HR & 95\% CrI & $\Prob(\mathrm{HR}{<}1)$ & $\Prob(\mathrm{HR}{<}0.8)$ \\
\midrule
Skeptical $\mathcal{N}(0,\,0.421^2)$          & $0$        & 0.610 & 0.485--0.757 & 1.000 & 0.994 \\
Weakly-inf.\ $\mathcal{N}(0,\,1)$       & $0$        & 0.591 & 0.466--0.738 & 1.000 & 0.997 \\
Enthusiastic $\mathcal{N}(\log 0.65,\,0.421^2)$ & $-0.431$ & 0.592 & 0.474--0.734 & 1.000 & 0.997 \\
Frequentist (adj.\ Cox, MI) [ref.] & -- & 0.621 & 0.496--0.776 & -- & -- \\
\bottomrule
\end{tabular*}
\end{table}

\begin{figure}[htbp]
  \centering
  \includegraphics[width=0.95\textwidth]{fig_bayes_priors.pdf}
  \caption{Prior sensitivity of the Bayesian treatment-effect
  estimate. Left: the three prior densities on the hazard-ratio
  scale. Right: the corresponding posterior densities, with the
  frequentist point estimate (dotted) for reference.}
  \label{fig:bayesprior}
\end{figure}

\subsubsection*{Bayesian sequential monitoring}

Figure~\ref{fig:bayesseq} replays evidence accumulation under the
skeptical prior at four event-driven looks. By the 25\% look (79
events), the posterior mean hazard ratio is already $0.601$ with
$\Prob(\mathrm{HR}<1) = 0.998$ and a credible interval excluding 1
(upper bound $0.88$); from the 50\% look (158 events) onward the
probability is essentially 1. Under the common Bayesian monitoring rule
of stopping when $\Prob(\mathrm{HR}<1)$ exceeds $0.975$, this trial
would have stopped at the \emph{first quarter} of the evidence---even
earlier than the frequentist group-sequential re-enactment of
Section~\ref{sec:res-design}, which required 50\% of the information and
a much larger $Z$. The contrast is partly definitional (the Bayesian
posterior probability does not pay a multiplicity price for looks) and
partly prior-driven (the skeptical prior pulls toward the null early
on, making its rapid crossing more impressive); both mechanisms deserve
emphasis when Bayesian designs are proposed as replacements for
alpha-spending.

\begin{figure}[htbp]
  \centering
  \includegraphics[width=0.95\textwidth]{fig_bayes_sequential.pdf}
  \caption{Bayesian sequential monitoring under the skeptical prior.
  Left: posterior mean hazard ratio with 95\% credible intervals at
  each look. Right: the monitoring statistic $\Prob(\mathrm{HR}<1)$
  against the 0.975 stopping threshold.}
  \label{fig:bayesseq}
\end{figure}

\subsubsection*{Model criticism}

The posterior predictive check (Figure~\ref{fig:bayesppc}) replicates
the control arm's Kaplan--Meier curve from 100 posterior draws: the
observed curve lies comfortably within the 95\% posterior predictive
band throughout follow-up, with no systematic deviation in the tail
where a poorly specified baseline hazard would typically betray itself.
Approximate leave-one-out cross-validation strongly prefers the Weibull
baseline over the exponential special case (ELPD difference $-18.0$, all
stacking weight on the Weibull), and the posterior shape parameter
$\alpha = 0.765$ (95\% CrI $0.695$--$0.837$) confirms a genuinely
decreasing baseline hazard---patients who survive the first two years
have progressively lower event rates, exactly as clinical experience
suggests.

\begin{figure}[htbp]
  \centering
  \includegraphics[width=0.90\textwidth]{fig_bayes_ppc.pdf}
  \caption{Posterior predictive check for the control arm: replicated
  Kaplan--Meier curves from 100 posterior draws versus the observed
  curve.}
  \label{fig:bayesppc}
\end{figure}

\subsection{Heterogeneous treatment effects}\label{sec:res-cate}

\subsubsection*{Average effects on the five-year absolute risk}

Table~\ref{tab:cate} reports the average treatment effect of levamisole
plus 5-FU versus observation on the five-year DFS event risk, estimated
three ways on the jackknife pseudo-observations ($n = 594$ complete
cases). The three estimators agree closely: the unadjusted
Kaplan--Meier difference is $-16.8$ percentage points (bootstrap 95\%
CI $-24.5$ to $-9.1$), cross-fitted G-computation gives $-15.3$ points,
and the doubly robust AIPW estimator gives $-16.2$ points with
influence-function standard error $3.9$ (95\% CI $-23.9$ to $-8.5$).
The consistency across a design-free estimator, an outcome-model-based
estimator, and a doubly robust estimator is exactly what one hopes to
see in a randomized comparison, and the magnitudes match the
intention-to-treat analysis of Section~\ref{sec:res-tx}.

\begin{table}[htbp]
\centering\small
\caption{Average treatment effect on the five-year DFS event risk
(negative $=$ risk reduction with levamisole plus 5-FU).}
\label{tab:cate}
\begin{tabular*}{\textwidth}{@{\extracolsep{\fill}} l c c c}
\toprule
Estimator & ATE (percentage points) & SE & 95\% CI \\
\midrule
Kaplan--Meier difference (bootstrap)   & $-16.8$ & 4.0 & $-24.5$ to $-9.1$ \\
G-computation (cross-fitted, pseudo-obs) & $-15.3$ & 0.5 & $-16.2$ to $-14.4$ \\
AIPW (cross-fitted, doubly robust)      & $\mathbf{-16.2}$ & 3.9 & $-23.9$ to $-8.5$ \\
\bottomrule
\end{tabular*}
\end{table}

\subsubsection*{Conditional effects: where the benefit lives}

The cross-fitted T-learner decomposes this average into patient-level
predicted benefits (Figure~\ref{fig:catedist}). The distribution is
wide: predicted five-year risk reductions range from negligible to over
40 percentage points, with quartile means of $28.9$, $18.6$, $12.6$,
and $1.1$ points. Crucially, the predicted benefit is strongly tied to
baseline risk (correlation $0.648$ between control-arm predicted risk
and predicted benefit): the patients predicted to benefit most are those
whose disease is most aggressive---high nodal burden, advanced local
extent---while low-risk patients have little room for absolute
improvement even under a proportionally constant treatment effect.
Table~\ref{tab:cateq} sharpens the picture with observed outcomes:
within the highest predicted-benefit quartile (mean baseline risk
$71.0\%$), the observed five-year risk difference is $-12.8$ points;
within the lowest quartile (baseline risk $46.6\%$), it is $-19.3$
points---statistically compatible with the same relative effect
everywhere. The apparent paradox (large predicted but modest observed
benefit in the top quartile, the reverse in the bottom) reflects both
estimation noise in quartile-specific contrasts (each based on
$\sim$64--98 events) and the mechanical relationship between baseline
risk and achievable absolute benefit.

\begin{figure}[htbp]
  \centering
  \includegraphics[width=0.95\textwidth]{fig_cate_distribution.pdf}
  \caption{Predicted absolute benefit of levamisole plus 5-FU at five
  years. Left: distribution across patients (AIPW average marked).
  Right: predicted benefit versus predicted baseline risk, coloured by
  randomized arm.}
  \label{fig:catedist}
\end{figure}

\begin{table}[htbp]
\centering\footnotesize
\caption{Predicted-benefit quartiles from the cross-fitted T-learner.
``Predicted benefit'' is the model-implied five-year risk reduction;
``observed risk difference'' is the corresponding pseudo-observation
contrast by arm within quartile.}
\label{tab:cateq}
\begin{tabular*}{\textwidth}{@{\extracolsep{\fill}} l c c c c c}
\toprule
Quartile & $n$ & Pred.\ benefit (pts) & Baseline risk & Observed\ diff (pts) & Events \\
\midrule
Q1 (largest predicted benefit) & 149 & 28.9 & 71.0\% & $-12.8$ & 98 \\
Q2                             & 148 & 18.6 & 58.2\% & $-16.1$ & 79 \\
Q3                             & 148 & 12.6 & 52.0\% & $-11.8$ & 72 \\
Q4 (smallest)                  & 149 & 1.1  & 46.6\% & $-19.3$ & 64 \\
\bottomrule
\end{tabular*}
\end{table}

The synthesis is the section's key message: in this trial, heterogeneity
in \emph{absolute} benefit is real and large, but it is driven almost
entirely by heterogeneity in \emph{baseline risk} rather than by
qualitative effect modification---consistent with the null interaction
tests of Section~\ref{sec:res-design}. For clinical decision-making,
this is the most useful kind of heterogeneity: it means a single
relative effect, applied to an accurate prognostic model, suffices to
identify who gains most in absolute terms, with no need to invoke
distinct treatment-response biology. All such CATE statements remain
exploratory in a single trial of this size.
